\begin{lemma}
Let $H\colon\mathbb{C}\to\mathbb{C}$ be real-differentiable at
$\xi\in\mathbb{C}$ with $\xi\neq0$, and write $r=|\xi|>0$ and
$e^{i\theta}=\xi/r$, so that $e^{-i\theta}=\bar\xi/r$.  Then
\[
  H_{\bar\xi}(\xi)=\frac12\Bigl(H_r(\xi)+\frac ir H_\theta(\xi)\Bigr)e^{i\theta},
  \qquad
  H_\xi(\xi)=\frac12\Bigl(H_r(\xi)-\frac ir H_\theta(\xi)\Bigr)e^{-i\theta},
\]
and equivalently
\[
  H_r(\xi)=H_\xi(\xi)\,e^{i\theta}+H_{\bar\xi}(\xi)\,e^{-i\theta},
\]
where $H_\xi$ is as in \noderef{DZ}, $H_{\bar\xi}$ as in \noderef{DZbar},
$H_r$ as in \noderef{HRad} and $H_\theta$ as in \noderef{HAng}.

The first two identities are the transformation formulas displayed at the start
of the proof of Theorem~2 in Section~4 of the paper.  The third is the
resolution of that pair for $H_r$; it is recorded here because the distortion
estimates (26), (27) and the Fourier computation of Section~4 use $H_r$
expressed through $H_\xi$ and $H_{\bar\xi}$ rather than the other way round.
\end{lemma}
\begin{proof}
Put $r=|\xi|$, which is $>0$ since $\xi\neq0$, and
\[
  e=\frac{\xi}{r},
\]
so that $|e|=1$, that is $e\bar e=1$, and $\bar e=\bar\xi/r$.  Write
$L=DH(\xi)$ for the real Fr\'echet derivative of $H$ at $\xi$, a real-linear
map $\mathbb{C}\to\mathbb{C}$, and put
\[
  A=H_\xi(\xi),\qquad B=H_{\bar\xi}(\xi).
\]
By \noderef{DZ} and \noderef{DZbar},
\begin{equation}\tag{1}
  A=\tfrac12\bigl(L[1]-i\,L[i]\bigr),\qquad B=\tfrac12\bigl(L[1]+i\,L[i]\bigr).
\end{equation}
Adding and subtracting the two formulas in (1) gives
\begin{equation}\tag{2}
  L[1]=A+B,\qquad L[i]=i(A-B),
\end{equation}
the second because $A-B=-i\,L[i]$ and $i(-i\,L[i])=L[i]$.

\emph{Step 1: the two directional derivatives.}  By \noderef{HRad},
$H_r(\xi)=L[\xi/|\xi|]=L[e]$, and by \noderef{HAng},
$H_\theta(\xi)=L[i\xi]=L[i\,re]=r\,L[ie]$, the last step by real-linearity of
$L$, $r$ being a real scalar.

Write $e=c+is$ with $c,s\in\mathbb{R}$, so that $ie=-s+ic$ and
$\bar e=c-is$.  By real-linearity of $L$ and (2),
\[
  H_r(\xi)=L[c\cdot1+s\cdot i]=c\,L[1]+s\,L[i]=c(A+B)+s\,i(A-B)
  =A(c+is)+B(c-is),
\]
that is
\begin{equation}\tag{3}
  H_r(\xi)=A\,e+B\,\bar e .
\end{equation}
Likewise
\[
  L[ie]=L[(-s)\cdot1+c\cdot i]=-s\,L[1]+c\,L[i]=-s(A+B)+c\,i(A-B)
  =A(-s+ic)+B(-s-ic),
\]
and $-s+ic=ie$ while $-s-ic=-i(c-is)=-i\bar e$, so $L[ie]=A(ie)-B(i\bar e)$ and
therefore
\begin{equation}\tag{4}
  H_\theta(\xi)=r\,L[ie]=i\,r\bigl(A\,e-B\,\bar e\bigr).
\end{equation}

\emph{Step 2: the two Wirtinger formulas.}  Since $r\neq0$, (4) gives
\[
  \frac ir H_\theta(\xi)=\frac ir\cdot i\,r\bigl(A\,e-B\,\bar e\bigr)
  =i^2\bigl(A\,e-B\,\bar e\bigr)=-A\,e+B\,\bar e .
\]
Adding and subtracting this and (3),
\[
  H_r(\xi)+\frac ir H_\theta(\xi)=2B\,\bar e,
  \qquad
  H_r(\xi)-\frac ir H_\theta(\xi)=2A\,e .
\]
Multiplying the first by $\tfrac12 e$ and the second by $\tfrac12\bar e$, and
using $e\bar e=1$,
\[
  \frac12\Bigl(H_r(\xi)+\frac ir H_\theta(\xi)\Bigr)e=B\,\bar e\,e=B=H_{\bar\xi}(\xi),
\]
\[
  \frac12\Bigl(H_r(\xi)-\frac ir H_\theta(\xi)\Bigr)\bar e=A\,e\,\bar e=A=H_\xi(\xi).
\]
Since $e=\xi/r=e^{i\theta}$ and $\bar e=\bar\xi/r=e^{-i\theta}$, these are the
first two identities of the statement, and (3) is the third.
\end{proof}
